% !TEX root = ../algebra_02.tex

\section{Инъективные и сюръективные линейные отображения. Принцип Дирихле для линейных операторов}

Ядро и образ служат точными мерами инъективности и сюръективности отображения соответственно.
Линейное отображение $\mathcal{A}$ называется мономорфизмом, если оно инъективно, и эпиморфизмом, если сюръективно.

\begin{Proposition}{Критерий инъективности}{}
	Линейное отображение $\mathcal{A} \in \mathrm{Hom}(V, W)$ инъективно $\;\Longleftrightarrow\; \Ker\mathcal{A} = \{0\}$.
\end{Proposition}
\begin{proof}
	$(\Rightarrow)$: При инъективности прообраз каждого элемента содержит не более одного вектора. Так как $\mathcal{A}(0) = 0$, то прообраз нуля состоит ровно из одного нулевого вектора, т.е. $\Ker\mathcal{A} = \{0\}$.
	
	$(\Leftarrow)$: Пусть $\mathcal{A}v_1 = \mathcal{A}v_2$. В силу линейности $\mathcal{A}(v_1 - v_2) = 0$, то есть $v_1 - v_2 \in \Ker\mathcal{A} = \{0\}$, откуда $v_1 = v_2$.
\end{proof}

Поведение базиса при отображении полностью характеризует его свойства:
\begin{Proposition}{}{}
	Пусть $(e_1, \ldots, e_n)$~--- базис $V$. Тогда:
	\begin{enumerate}
		\item $\mathcal{A}$ инъективно $\;\Leftrightarrow\;$ векторы $\mathcal{A}e_1, \ldots, \mathcal{A}e_n$ \textbf{линейно независимы} в $W$.
		\item $\mathcal{A}$ сюръективно $\;\Leftrightarrow\;$ векторы $\mathcal{A}e_1, \ldots, \mathcal{A}e_n$ являются \textbf{системой образующих} для $W$.
		\item $\mathcal{A}$ биективно (изоморфизм) $\;\Leftrightarrow\;$ векторы $\mathcal{A}e_1, \ldots, \mathcal{A}e_n$ образуют \textbf{базис} $W$.
	\end{enumerate}
\end{Proposition}

\begin{proof}
	\begin{enumerate}
		\item
		\begin{itemize}
			\item $(\Rightarrow)$: Пусть $\lambda_1\,\mathcal{A}e_1 + \ldots + \lambda_n\,\mathcal{A}e_n = 0$.
			Тогда $\mathcal{A}(\lambda_1 e_1 + \ldots + \lambda_n e_n) = 0$,
			и из $\Ker\mathcal{A} = \{0\}$ следует $\lambda_1 e_1 + \ldots + \lambda_n e_n = 0$,
			откуда $\lambda_1 = \ldots = \lambda_n = 0$.
			\item $(\Leftarrow)$: Пусть $\mathcal{A}v = 0$, $v = \lambda_1 e_1 + \ldots + \lambda_n e_n$.
			Тогда $\lambda_1\,\mathcal{A}e_1 + \ldots + \lambda_n\,\mathcal{A}e_n = 0$,
			и из линейной независимости $\lambda_1 = \ldots = \lambda_n = 0$, то есть $v = 0$.
		\end{itemize}
		\item
		\begin{itemize}
			\item $(\Rightarrow)$: Для любого $w \in W$ найдём $v \in V$ с $\mathcal{A}v = w$.
			Разложив $v$ по базису:
			\[
			w = \mathcal{A}(\lambda_1 e_1 + \ldots + \lambda_n e_n) = \lambda_1\,\mathcal{A}e_1 + \ldots + \lambda_n\,\mathcal{A}e_n.
			\]
			\item $(\Leftarrow)$: Для любого $w \in W$ запишем $w = \lambda_1\,\mathcal{A}e_1 + \ldots + \lambda_n\,\mathcal{A}e_n = \mathcal{A}(\lambda_1 e_1 + \ldots + \lambda_n e_n) \in \Im\mathcal{A}$.
		\end{itemize}
		\item Следует из пунктов 1 и 2.
	\end{enumerate}
\end{proof}

\begin{Statement}{Принцип Дирихле для линейных операторов}{}
	Пусть размерности пространств совпадают и конечны: $\dim V = \dim W = n < \infty$, и $\mathcal{A} \in \mathrm{Hom}(V, W)$.
	Тогда следующие три условия эквивалентны:
	\begin{enumerate}
		\item $\mathcal{A}$~--- инъективно (мономорфизм),
		\item $\mathcal{A}$~--- сюръективно (эпиморфизм),
		\item $\mathcal{A}$~--- биективно (изоморфизм).
	\end{enumerate}
\end{Statement}
\begin{proof}
	Согласно теореме о ранге и дефекте:
	\[
	\dim\Ker\mathcal{A} + \dim\Im\mathcal{A} = n.
	\]
	Инъективность означает $\Ker\mathcal{A} = \{0\}$, то есть $\dim\Ker\mathcal{A} = 0$.
	Подставляя это в формулу, получаем $\dim\Im\mathcal{A} = n$. Поскольку $\dim W = n$, это равносильно тому, что $\Im\mathcal{A} = W$, то есть отображение сюръективно.
	Следовательно, любое инъективное или сюръективное отображение между пространствами равной размерности автоматически является биективным.
\end{proof}
